\documentclass[10pt]{amsart}
\usepackage[a4paper,margin=23mm]{geometry}
\usepackage{amsmath,amssymb,amsthm,mathtools,hyperref,graphicx,xspace,letltxmacro,xpatch}
\usepackage[english]{babel}
\newtheorem{theo}{Theorem}[section]
\newtheorem{lemm}[theo]{Lemma}\newtheorem{prop}[theo]{Proposition}\newtheorem{coro}[theo]{Corollary}\newtheorem{conj}[theo]{Conjecture}
\theoremstyle{definition}\newtheorem{defi}[theo]{Definition}\newtheorem{rema}[theo]{Remark}\newtheorem{exam}[theo]{Example}\newtheorem{exem}[theo]{Example}\newtheorem{nota}[theo]{Notation}
\newcommand{\ie}{i.e.}\newcommand{\eg}{e.g.}\newcommand{\resp}{resp.}\newcommand{\skpt}{\par\smallskip\noindent}
\emergencystretch=3em\pagestyle{plain}



\newcommand{\bbP}{\mathbb{P}}
\newcommand{\bbR}{\mathbb{R}}
\newcommand{\bbC}{\mathbb{C}}
\newcommand{\bbN}{\mathbb{N}}

\newcommand{\clC}{\mathcal{C}}

\newcommand{\clF}{\mathcal{F}}
\newcommand{\frkm}{\mathfrak{m}}

\newcommand{\bfun}{\mathbf{1}}

\newcommand{\noopsort}[1]{}

%%%%%%%%%%%%%%%%%%%%%%%%%%%%%%%%%%%%%%%%%%%%%%%%%%%%%%%%%%%%%%%%%%%%

\graphicspath{{./figures/}}

\newcommand*{\mk}{\mkern -1mu}
\newcommand*{\Mk}{\mkern -2mu}
\newcommand*{\mK}{\mkern 1mu}
\newcommand*{\MK}{\mkern 2mu}

%\hypersetup{urlcolor=purple, linkcolor=blue, citecolor=red}

\let\oldhat\hat
\renewcommand*{\hat}[1]{\mathchoice{\widehat{#1}}{\widehat{#1}}{\oldhat{#1}}{\oldhat{#1}}}
\let\oldforall\forall
\renewcommand*{\forall}{\mathrel{\oldforall}}

%%%%%%%%%%%%%%%%%%%%%%%%%%%%%%%%%%%%%%%%%%%%%%%%%%%%%%%%%%%%%%%%%%%%

\title{Global L2 Cauchy theorem for nonlinear Schrodinger equations with fractional-Brownian dispersion}
\author{Duboscq, Romain, Réveillac, Anthony}
\date{doi:10.5802/ahl.122; CC BY 4.0}
\begin{document}
\maketitle
\noindent\textit{Editorial note.} Only selected author theorem counted; complete original mathematical body and bibliography preserved. No strengthening of original hypotheses or conclusions.\par\medskip
\noindent\textit{Allowed standard prerequisites (editorial).}\begin{itemize}
\item Standard functional analysis, measure theory and the author-cited classical theorems. Standard prior results stay cited; complete source-local proofs and definitions retained.
\end{itemize}\medskip

\section{Introduction}

Let $(\Omega, \bbP)$ be the standard probability space endowed with the Wiener filtration $(\clF_t)_{t\,\geq\,0}$. We consider the stochastic process $W^H$, a fractional Brownian motion with Hurst parameter $H\in(0,1)$, given by, $\forall t\in\bbR^+$,
\[
W^H_t = \int_{-\infty}^t\left((t-s)_+^{H}-(-s)_+^{H}\right)dW_s,
\]
where $W$ is a standard wiener process. The main objective of this paper is to derive a stochastic counter-part to the classical Hardy--Littlewood--Sobolev inequality~\cite{hardy1928some,hardy1932some,sobolev1938theorem}. To be more specific, we obtain the following result.

\begin{theo}\label{thm:HLSsto}
Let $(W^H_t)_{t\,\geq\,0}$ be a fractional Brownian motion of Hurst index $H\in (0,1)$, $\beta\in (0,1-H)$, $p,q \in (1,\infty)$ and $\alpha\in (0,1)$ such that
\[
2 - \alpha = \frac1p+\frac1q.
\]
Then, there exist $T>0$ and $C_{\ref{thm:HLSsto}}>0$ such that, $\bbP$-a.s., $\forall f\in L^p([0,T])$, $\forall g\in L^q([0,T])$ the following inequality holds
\begin{equation}\label{ineq:HLSsto}
\left|\int_0^T\int_0^T f(t)\left|W^{H}_t-W^{H}_s\right|^{-\alpha}g(s) ds dt\right|\leq C_{\ref{thm:HLSsto}} T^{\alpha\beta}\|f\|_{L^p\left([0,\,T]\right)}\|g\|_{L^q\left([0,\,T]\right)}.
\end{equation}
\end{theo}

Our motivation to prove such result stems from the Cauchy problem of nonlinear evolution equations with a randomly modulated dispersion. Such equations are for instance: the nonlinear Schr\"odinger equation~\cite{belaouar2015numerical,chouk2015nonlinear,de2010nonlinear,debussche20111d}, the Korteweg--de Vries equation~\cite{chouk2015nonlinear} and the Benjamin--Bona--Mahony~\cite{chen2017generalized}. They have recently raised an interest due to the effects of the stochastic modulation. Here, we address the local Cauchy problem for the following nonlinear Schr\"odinger equation with noisy dispersion in its mild formulation
\begin{equation}\label{eq:main}
\psi(t,x) = P_{0,\,t} \psi_0(x) + \lambda \int_0^t P_{s,\,t} |\psi|^{2\sigma} \psi(s,x)ds,\quad \forall (t,x)\in[0,T]\times\bbR^d,
\end{equation}
where $\lambda \in\bbC$ and, $\forall \varphi\in\clC^{\infty}_0(\bbR^d)$,
\[
P_{s,\,t}\varphi : = \clF^{-1}\left(e^{-i|\xi|^2\left(W^H_t-W^H_s\right)} \hat{\varphi}\right).
\]

For $d=1$, $H=1/2$ and $\sigma = 1$, this equation arises in the field of nonlinear optics as a limit model for the propagation of light pulse in an optical fiber where the dispersion varies along the fiber~\cite{agrawal2001applications,agrawal2007nonlinear}. These variations in the dispersion accounts for the so-called \emph{dispersion management} which aims to improve the transmission of a light signal by constructing a zero-mean dispersion fiber in order to avoid the problem of the chromatic dispersion of the light signal. When the variations are assumed to be random, a noisy dispersion can be derived (see~\cite{de2010nonlinear,marty2006splitting}) which leads, in the white noise case, to Equation~\eqref{eq:main}.

As part of the problems concerning the propagation of waves in random media, there is a vast literature around random Schr\"odinger equations. Let us mention in particular the cases of random potentials~\cite{erdHos2008quantum,erdHos2000linear} and noisy potentials~\cite{de2002effect,de2005blow,de1999stochastic}. In these works, the effects of the stochastic potential greatly affect the dynamic of the Schr\"odinger equation and are, in a broader context, a motivation to introduce randomness in PDEs. Specifically, there is a well known effect which attracted a lot of attention: the so-called \emph{regularization by noise} phenomenon (see~\cite{flandoli2011random} for a survey). This phenomenon can be summarized as an improvement, due to the presence of noise, of the well-posedness of differential equations and has been studied in the context of SDEs
~\cite{CATELLIER20162323,krylov2005strong,priola2012pathwise,veretennikov1981strong,
zvonkin1974transformation}, transport equation~\cite{catellier2016rough,fedrizzi2013noise,flandoli2010well}, SPDEs~\cite{da2013strong} and scalar conservation laws~\cite{gess2017long}. We remark that obtaining a regularization by noise in the context of nonlinear random PDE is a challenging task and most of the results are obtained in a linear setting. For instance, an open problem is to obtain a regularization by noise for the Euler or Navier--Stokes equations.

As mentioned previously, we are not the first one to investigate the Cauchy problem of Equation~\eqref{eq:main}. It was first studied in~\cite{de2010nonlinear} where the global Cauchy problem was solved for $H = 1/2$ and $\sigma < 2/d$ (which corresponds to a classical $L^2$-subcritical nonlinearity). In~\cite{debussche20111d}, the authors proved that, in the $L^2$-critical case, when $d = 1$, $H = 1/2$ and $\sigma = 5$, the solutions are globally well-posed, which is not the case for the deterministic nonlinear Schr\"odinger equation and, thus, hints for a regularization by noise effect. In~\cite{chouk2015nonlinear}, the authors study the case for $d=1$, $\sigma = 2$, $H$ small enough and $d=2$, $\sigma = 1$, $H\in(0,1)$. By a simple scaling argument on the space and time variables of~\eqref{eq:main} and thanks to the scaling invariance of the Wiener process, it was conjectured in~\cite{belaouar2015numerical} that, in fact, the critical nonlinearity should be $\sigma = 4/d$ for $H= 1/2$, a $L^2$-supercritical nonlinearity, which is twice as large as the deterministic $L^2$-critical nonlinearity. Furthermore, this fact was supported by numerical simulations in $1$D and leads to believe that the white noise dispersion has a strong stabilizing property.

In this paper, we prove the global Cauchy problem~\eqref{eq:main} for $d\in\bbN$, $\sigma\leq 2/d$ and $H\in(0,1)$. To be more specific, we obtain the following result.


\begin{theo}\label{thm:main}
Let $\sigma\leq \frac{2}{d}$, $\psi_0\in L^2(\bbR^d)$ and $a\in (2,\infty)$ such that $2/a = d\left(1/2-1/(2\sigma+2)\right)$. Then, $\bbP$-a.s., there exists a unique solution $\psi$ which belongs in $L^{a}([0,+\infty[;L^{2\sigma+2}(\bbR^d))\,\cap\, \clC([0;+\infty[;L^2(\bbR^d))$ to~\eqref{eq:main}.
\end{theo}


\begin{rema}
Thus, the modulation by a random noise of the dispersion operator leads to a regularizing effect since we are able to prove the global existence and uniqueness of solutions in the critical case $\sigma = 2/d$ for arbitrarily large initial data $\psi_0\in L^2(\bbR^d)$. The other interesting fact of our result is that, no matter how close to $1$ the Hurst parameter is, we still reach the critical case. This problem was left open in~\cite{chouk2015nonlinear} where $H$ needs to be small enough.
\end{rema}


The classical approach to investigate the Cauchy problem for nonlinear Schr\"odinger equations is to derive local Strichartz estimates~\cite{cazenave2003semilinear}. These estimates are a direct consequence of the dispersive property of the linear operator $i\Delta$. However, as pointed out in~\cite{debussche20111d}, it is much harder to obtain such estimates in the case of a white noise dispersion because of the presence of the Wiener process. We recall from~\cite{de2010nonlinear,marty2006splitting} that the propagator associated to the linear part of~\eqref{eq:main} is explicitly given by, $\forall t,s\in (0,\infty)$ and $\forall \varphi\in\clC^{\infty}_0(\bbR^d)$,
\begin{equation}\label{eq:propagator}
P_{s,\,t}\varphi(x) = \frac{1}{\left(4\pi\left(W_t^H-W_s^H\right)\right)^{d/2}}\int_{\bbR^d} e^{i\frac{|x-y|^2}{4\left(W_t^H-W_s^H\right)}} \varphi(y)dy.
\end{equation}

Following the classical proof of Strichartz estimates (see for instance~\cite{keel1998endpoint}), a fundamental tool is the Hardy--Littlewood--Sobolev inequality. This is where Theorem~\ref{thm:HLSsto} comes at hand since the classical potential $|t-s|^{-\alpha}$ is replaced by $|W_t-W_s|^{-\alpha}$. As a direct consequence, we obtain the following stochastic Strichartz estimates.

\begin{defi}
For any $(q,p)\in (2,\infty)^2$, we say that $(q,p)$ is an admissible pair if
\[
\frac{2}{q} = d\left(\frac 1 2 - \frac 1 p \right).
\]
\end{defi}

\begin{prop}\label{prop:Strichartz}
Let $(W^H_t)_{t\geq 0}$ be a fractional Brownian motion of Hurst index $H\in (0,1)$, $\beta\in (0,1-H)$, $(q,p)$ an admissible pair and $\alpha = d(1/2-1/p)=2/q$. Then, $\bbP$-a.s., there exist $T>0$ and $C_{1,\,\ref{prop:Strichartz}},C_{2,\,\ref{prop:Strichartz}}>0$ such that, $\forall f\in L^2(\bbR^d)$ and $\forall g\in L^{r'}([0,T];L^{l'}(\bbR^d))$, the following inequalities holds
\begin{align}
\left\|P_{0,\,\cdot} f \right\|_{L^q\left([0,\,T]\,;\,L^p\left(\bbR^d\right)\right)}&\leq C_{1,\,\ref{prop:Strichartz}}T^{\alpha\beta}\|f\|_{L^2}\label{eq:Stri1}
,
\\
\left\|\int_0^T P_{s,\,\cdot}g(s)ds \right\|_{L^q\left([0,\,T]\,;\,L^p\left(\bbR^d\right)\right)} &\leq C_{2,\,\ref{prop:Strichartz}}T^{\alpha\beta} \|g\|_{L^{r'}\left([0,\,T]\,;\,L^{l'}\left(\bbR^d\right)\right)},\label{eq:Stri2}
\end{align}
for any $(r,l)$ admissible pair.
\end{prop}

These Strichartz estimates are more powerful due to the presence of the term~$T^{\alpha\beta}$. Indeed, in the fixed-point argument, this term will be necessary to obtain the contraction of the mapping in the critical case $\sigma = 2/d$ and, as opposed to the deterministic case, it will only depend on the size of the initial data (even in the critical case). This will be the main argument to prove Theorem~\ref{thm:main}. The rest of the paper is devoted to the proof of Theorem~\ref{thm:HLSsto} in Section~\ref{sec:proofThmHLSsto} and the proofs of Proposition~\ref{prop:Strichartz} and Theorem~\ref{thm:main} in Section~\ref{sec:proofThmMain}.

\section{Proof of Theorem~\ref{thm:HLSsto}}\label{sec:proofThmHLSsto}

Before proceeding any further, let us remark that we can, without loss of generality, assume that $f\in L^p([0,T])$ and $g\in L^q([0,T])$ are non-negative functions and, furthermore, by a scaling argument, we can assume that $\|f\|_{L^p([0,\,T])} = \|g\|_{L^q([0,\,T])} = 1$. Our strategy follows the proof of Lieb--Loss~\cite[Section~4.3]{lieb2001analysis}. It is based on the following layer cake representation
\begin{align*}
f(t) &= \int_0^{+\infty} \bfun_{f(t)\,>\,a}\,da,\quad g(s) = \int_0^{+\infty} \bfun_{g(s)\,>\,b}\,db
\\ \text{and}\quad \left|W_t^H-W_s^H\right|^{-\alpha} &= \alpha \int_0^{+\infty} c^{-1-\alpha} \bfun_{\left|W_t^H-W_s^H\right|\,<\,c}\,dc.
\end{align*}
Then, by Fubini's theorem, the left-hand-side of~\eqref{ineq:HLSsto} can be rewritten as
\begin{multline}\label{eq:ineqHLScake}
\int_0^T\int_0^T f(t)\left|W_t^H-W_s^H\right|^{-\alpha}g(s) ds dt =\\
\alpha\int_0^{+\infty}\int_0^{+\infty}\int_0^{+\infty} c^{-1-\alpha} \left(\int_0^T\int_0^T \bfun_{f(t)\,>\,a}\bfun_{g(s)\,>\,b}\bfun_{\left|W_t^H-W_s^H\right|\,<\,c} ds dt\right) dc db da.
\end{multline}
By denoting
\begin{align*}
\check{f}(a) :=& \int_0^T \bfun_{f(t)\,>\,a}\,dt,\quad\check{g}(b) := \int_0^T \bfun_{g(s)\,>\,b}\,ds\\
\text{and}\quad \check{W}(c,T) :=& \sup_{t\,\in\,[0,\,T]} \int_0^T \bfun_{\left|W_t^H-W_s^H\right|\,<\,c}\,ds,
\end{align*}
we have the following result whose proof is postponed.

\begin{lemm}\label{lem:estWcheck}
Let $H\in(0,1)$ and $\beta\in (0,1-H)$. There exists $T>0$ and $C_{\ref{lem:estWcheck}}>0$ such that, $\bbP$-a.s., $\forall c\in\bbR^{+*}$,
\[
\check{W}(c,T) \leq C_{\ref{lem:estWcheck}} T^{\beta} c.
\]
\end{lemm}

We now set $p,q\in (2,+\infty)$ such that
\[
\frac 1 p + \frac 1 q + \alpha = 2.
\]
We see that we can bound each characteristic function by $1$ in~\eqref{eq:ineqHLScake} and, thus, we deduce that
\begin{align*}
\int_0^T\int_0^T f(t)\left|W_t^H-W_s^H\right|^{-\alpha}g(s) ds dt \leq \alpha\int_0^{+\infty}\int_0^{+\infty}\int_0^{+\infty} c^{-\alpha-1}I(a,b,c) da db dc,
\end{align*}
with
\[
I(a,b,c) : =
\begin{cases}
\check{f}(a)\check{g}(b), &\text{if}\quad \frkm(a,b,c) = \check{h}(c),
\\ \check{f}(a)\check{W}(c,T), &\text{if}\quad \frkm(a,b,c) = \check{g}(b),
\\ \check{W}(c,T)\check{g}(b), &\text{if}\quad \frkm(a,b,c) = \check{f}(a),
\end{cases}
\]
where
\[
\frkm(a,b,c) = \lim_{\iota\,\to\,0} \max\left\{ \check{f}(a),\check{g}(b),\check{h}(c)+\iota\right\}\quad\text{and}\quad\check{h}(c) : = C_{\ref{lem:estWcheck}} T^{\beta} c.
\]

\begin{rema}
In the previous definition of $\frkm$, we choose to have, in the case where $\max\{ \check{f}(a),\check{g}(b),\check{h}(c)\} = \check{f}(a) = \check{h}(c)$, $\frkm(a,b,c) = \check{h}(c)$ (and similarly for $\check{g}$).
\end{rema}

From here, we essentially follow the arguments from~\cite{lieb2001analysis}. We first assume that $\check{f}(a) \geq \check{g}(b)$. We deduce that
\begin{multline*}
\int_0^{+\infty} c^{-\alpha-1} I(a,b,c)dc \leq\\
\check{g}(b) \int_0^{+\infty} c^{-\alpha-1} \check{h}(c)\bfun_{\check{h}(c)\,\leq\,\check{f}(a)}dc + \check{f}(a)\check{g}(b)\int_0^{+\infty}c^{-\alpha-1}\bfun_{\check{f}(a)\,\leq\,\check{h}(c) }dc.
\end{multline*}
We denote $\kappa_{T} = C_{\ref{lem:estWcheck}} T^{\beta}$. Since $\check{h}(c) \leq \check{f}(a)$ is equivalent to
\begin{equation*}
c\leq \check{f}(a)/\kappa_{T},
\end{equation*}
the first integral on the right-hand-side is estimated as
\begin{align*}
\int_0^{+\infty} c^{-\alpha-1} \check{h}(c)\bfun_{\check{h}(c)\,\leq\, \check{f}(a)}dc&\leq \kappa_T \int_0^{\check{f}(a)/\kappa_{T}} c^{-\alpha}dc
\\ &\leq (1-\alpha)^{-1}\kappa_T\left(\check{f}(a)/\kappa_{T}\right)^{1 - \alpha} = (1-\alpha)^{-1}\kappa_T^{\alpha} \check{f}(a)^{1-\alpha}.
\end{align*}
The second integral is bounded as
\begin{align*}
\int_0^{+\infty}c^{-\alpha-1}\bfun_{\check{f}(a)\,\leq\,\check{h}(c) }dc \leq \int_{\check{f}(a)/\kappa_{T}}^{+\infty}c^{-\alpha-1}dc \leq \alpha^{-1}\kappa_T^{\alpha} \check{f}(a)^{-\alpha}.
\end{align*}
Hence, since by assumption $\check{g}(b)^{-\alpha}\geq \check{f}(a)^{-\alpha}$, it follows that
\[
\int_0^{+\infty} c^{-\alpha-1} I(a,b,c)dc \lesssim_{\alpha} \kappa_T^{\alpha} \check{g}(b)\check{f}(a)^{1-\alpha} = \kappa_T^{\alpha} \min\left\{\check{g}(b)\check{f}(a)^{1-\alpha},\check{f}(a)\check{g}(b)^{1-\alpha}\right\}.
\]
By assuming that $\check{f}(a) \leq \check{g}(b)$ and following the same arguments, we obtain
\[
\int_0^{+\infty} c^{-\alpha-1} I(a,b,c)dc \lesssim_{\alpha} \kappa_T^{\alpha} \min\left\{\check{g}(b)\check{f}(a)^{1-\alpha}, \check{f}(a)\check{g}(b)^{1-\alpha}\right\}.
\]
We then proceed by integrating with respect to $b$ and $a$. We have
\begin{multline*}
\int_0^{+\infty} \int_0^{+\infty} \min\left\{ \check{g}(b)\check{f}(a)^{1-\alpha}, \check{f}(a)\check{g}(b)^{1-\alpha}\right\} db da =\\
\int_0^{+\infty}\int_0^{a^{p/q}}\check{f}(a)\check{g}(b)^{1-\alpha}dbda+\int_0^{+\infty}\int_{a^{p/q}}^{+\infty} \check{g}(b)\check{f}(a)^{1-\alpha} db da.
\end{multline*}
Thanks to H\"older' inequality, we have, for $r = (q-1)(1-\alpha)$,
\begin{align*}
\int_0^{a^{p/q}}\check{g}(b)^{1-\alpha}db &= \int_0^{a^{p/q}}\check{g}(b)^{1-\alpha}b^{-r}b^{r}db
\\ &\leq \left(\int_0^{a^{p/q}}\check{g}(b) b^{q-1}db \right)^{1-\alpha}\left(\int_0^{a^{p/q}} b^{-r/\alpha}db\right)^{\alpha}.
\end{align*}
The norms of $f$ and $g$ are such that
\[
\|f\|_{L^p\left([0,\,T]\right)}^p = p\int_0^{+\infty} a^{p-1} \check{f}(a)da = 1\quad\text{and}\quad \|g\|_{L^q\left([0,\,T]\right)}^q = q\int_0^{+\infty} b^{q-1} \check{g}(b)db = 1.
\]
Thus, since
\begin{align*}
\frac pq\left(1-\frac{r}{\alpha}\right)\alpha &= \frac pq\big(\alpha - (q-1)\left(1-\alpha\right)\big)
\\ &=\frac pq\big(1 - q\left(1-\alpha\right)\big) = p\left(\frac1q - 1 + \alpha\right) = p - 1,
\end{align*}
we obtain
\begin{multline*}
\int_0^{+\infty}\check{f}(a)\int_0^{a^{p/q}}\check{g}(b)^{1-\alpha}db da\\ \lesssim \|g\|_{L^q\left([0,\,T]\right)}^{q(1-\alpha)}\int_{0}^{+\infty}\check{f}(a)a^{p-1}da = \|g\|_{L^q\left([0,\,T]\right)}^{q(1-\alpha)}\|f\|_{L^p\left([0,\,T]\right)}^p\leq 1.
\end{multline*}
By similar arguments, we deduce that
\[
\int_0^{+\infty}\int_{a^{p/q}}^{+\infty} \check{g}(b)\check{f}(a)^{1-\alpha} db da \lesssim 1,
\]
which concludes the proof of Theorem~\ref{thm:HLSsto}.


It remains to prove Lemma~\ref{lem:estWcheck}. We have that, $\forall c>0$ and $\forall t\in [0,T]$,
\[
\int_0^T \bfun_{\left|W_t^H-W_s^H\right|\,<\,c}\,ds = \int_{\bbR} \bfun_{|y|\,<\,c} \ell^{W_t^H-y}_{[0,\,T]}\,dy = \int_{-c}^{c} \ell^{W_t^H-y}_{[0,\,T]}\,dy,
\]
where $\ell$ is the local time of $W^H$ given as
\[
\ell^z_{[s,\,t]} : = \lim_{\varepsilon\,\to\,0}\int_{[s,\,t]} P_{\varepsilon} \delta_z\left(W^H_u\right)du,
\]
where $(P_{t})_{t\,\geq\,0}$ is the heat semigroup. We need the following result from~\cite{xiao1997holder}.

\begin{theo}\label{thm:LTHolder}
Let $(W^H_t)_{t\,\geq\,0}$ be a fractional Brownian motion with Hurst parameter $H\in(0,1)$. Then, for any interval $I\subset \bbR$, there exists a positive finite constant $K$ such that, $\bbP$-a.s.,
\[
\underset{r\,\to\,0}{\limsup} \sup_{t\,\in\,I} \frac{\sup_{x\,\in\,\bbR}\ell^x_{\left[t-r,\,t+r\right]}}{r^{1-H}\log(1/r)^H}\leq K.
\]
\end{theo}

We deduce from the previous result that for any $\beta\in (0,1-H)$ and $T>0$ small enough, there exists a constant $C_{\ref{thm:LTHolder}}>0$ such that, $\bbP$-a.s., we have
\[
\sup_{x\,\in\,\bbR} \ell^x_{[0,\,T]} \leq C_{\ref{thm:LTHolder}} T^{\beta}.
\]
It follows that, for $T$ small enough,
\[
\int_0^T \bfun_{\left|W_t^H-W_s^H\right|\,<\,c}ds \leq 2 C_{\ref{thm:LTHolder}} c T^{\beta},
\]
which is exactly the desired result.


\section{Proof of Proposition~\ref{prop:Strichartz} and Theorem~\ref{thm:main}}\label{sec:proofThmMain}

\subsection{Stochastic Strichartz estimates}

Since $(P_{s,\,t})_{0\,\leq\,s\,\leq\,t}$ is an isometry from $L^2$ to itself, we deduce by the Hausdorff--Young inequality and an interpolation argument, that, $\forall p\in [2,\infty]$,
\begin{equation}\label{ineq:dispest}
\left\|P_{s,\,t}\varphi\right\|_{L^p\left(\bbR^d\right)} \lesssim \frac{1}{\left|W_t-W_s\right|^{d\left(1/2-1/p\right)}} \|\varphi\|_{L^{p'}\left(\bbR^d\right)},
\end{equation}
where $p'$ is the H\"older conjugate of $p$. We denote $(P_{s,\,t}^*)_{0\,\leq\,s\,\leq\,t}$ the adjoint of $(P_{s,\,t})_{0\,\leq\,s\,\leq\,t}$, that is
\[
P^*_{s,\,t}\varphi(x) : = \clF^{-1}\left(e^{-\frac{i}2 |\xi|^2\left(W_t^H-W_s^H\right)} \hat{\varphi}(\xi)\right) = P_{t,\,s}.
\]
This leads, in particular, to the fact that
\[
P^*_{s,\,t} = P_{t,\,s},\quad P_{0,\,s}^*P_{0,\,t} = P_{s,\,t}\quad\text{and}\quad P_{s,\,t}P^*_{r,\,t} = P_{s,\,r},\; \forall r\in[s,t].
\]
The proof of Proposition~\ref{prop:Strichartz} is based on the $TT^*$ argument. We set $\alpha = d(1/2-1/p)$ and consider the integral, $\forall f,g\in \clC([0,T], \clC^{\infty}_0(\bbR^d))$,
\begin{align*}
I(f,g) : = \left| \int_0^T\int_0^T \big\langle P_{0,\,t} f(s), P_{0,\,s} g(t)\big\rangle_{L^2} ds dt \right| &= \left| \int_0^T\int_0^T \left\langle P_{0,\,s}^* f(s), P_{0,\,t}^*g(t)\right\rangle_{L^2} ds dt \right|
\\
&= \left| \int_0^T\int_0^T \big\langle P_{s,\,t} f(s), g(t)\big\rangle_{L^2} ds dt \right|
\end{align*}
By H\"older's inequality, \eqref{ineq:dispest} and Theorem~\ref{thm:HLSsto}, we deduce that
\begin{align*}
I(f,g) &\leq \int_0^T\int_0^T \left\|P_{s,\,t}f(s)\right\|_{L^{p}\left(\bbR^d\right)}\|g(t)\|_{L^{p'}\left(\bbR^d\right)} ds dt\\
&\lesssim\int_0^T\int_0^T \left|W_t-W_s\right|^{-\alpha}\|f(t)\|_{L^{p'}\left(\bbR^d\right)}\|g(s)\|_{L^{p'}\left(\bbR^d\right)} ds dt\\
&\lesssim T^{\alpha\beta} \|f\|_{L^{q'}\left([0,\,T],\,L^{p'}\left(\bbR^d\right)\right)}\|g\|_{L^{q'}\left([0,\,T],\,L^{p'}\left(\bbR^d\right)\right)},
\end{align*}
since
\[
2 - d\left(\frac12-\frac1p\right) = \frac 1{q'} + \frac1{q'} = 2 - \frac{2}{q}.
\]
This yields, on one hand, that
\begin{equation}\label{eq:StriHomoDual}
\left\|\int_0^T P_{0,\,s}^* f(s) ds \right\|_{L^2\left(\bbR^d\right)}^2 = I(f,f) \lesssim T^{\alpha\beta} \|f\|_{L^{q'}\left([0,\,T],\,L^{p'}\left(\bbR^d\right)\right)}^2,
\end{equation}
and, on another hand, by a duality argument,
\begin{equation}\label{eq:Stri2b}
\left\| \int_0^T P_{s,\,\cdot} f(s) ds \right\|_{L^{q}\left([0,T],\,L^p\left(\bbR^d\right)\right)} \lesssim T^{\alpha\beta}\|f\|_{L^{q'}\left([0,\,T],\,L^{p'}\left(\bbR^d\right)\right)}
\end{equation}
We are now in position to prove~\eqref{eq:Stri1} and~\eqref{eq:Stri2}. It follows from~\eqref{eq:StriHomoDual} that, $\forall f\in L^2(\bbR^d)$ and $\forall g\in L^{q'}([0,T];L^{p'}(\bbR^d))$,
\begin{align*}
\int_{0}^T \left\langle P_{0,\,t}f,g(t) \right\rangle_{L^2} dt &= \left\langle f, \int_0^T P_{0,\,t}^*g(t)\right\rangle_{L^2} \leq \|f\|_{L^2(\bbR^d)} \left\|\int_0^T P_{0,\,t}^* g(t) ds \right\|_{L^2\left(\bbR^d\right)}^2\\
&\lesssim T^{\alpha\beta}\|f\|_{L^2\left(\bbR^d\right)} \|g\|_{L^{q'}\left([0,\,T],\,L^{p'}\left(\bbR^d\right)\right)},
\end{align*}
we obtain~\eqref{eq:Stri1} by a duality argument. We now turn to~\eqref{eq:Stri2}. We have, by~\eqref{eq:StriHomoDual},
\begin{align*}
\left\|\int_0^T P_{s,\,\cdot} f(s) ds\right\|_{L^q\left([0,\,T]\,;\,L^p\left(\bbR^d\right)\right)} &\leq \int_0^T \left\|P_{s,\,\cdot} f(s) \right\|_{L^q\left([0,\,T]\,;\,L^p\left(\bbR^d\right)\right)} ds
\\
&\lesssim T^{\alpha\beta} \int_0^T \|f(s)\|_{L^2\left(\bbR^d\right)} ds = T^{\alpha\beta} \|f\|_{L^1\left([0,\,T]\,;\,L^2\left(\bbR^d\right)\right)}.
\end{align*}
Thanks to this estimate, by an interpolation argument with~\eqref{eq:Stri2b}, we deduce~\eqref{eq:Stri2}.

\subsection{Well-posedness of equation~\eqref{eq:main}}

We can now apply the previous result to solve the local Cauchy problem of~\eqref{eq:main} The strategy is based on a fixed-point argument of the mapping $\Gamma$ from $L^{q}([0,T];L^p(\bbR^d))$ to itself given by
\begin{equation}\label{eq:GammaMap}
\Gamma(\psi)(t,x) = P_{0,\,t}\psi_0(x) - i\lambda \int_0^T P_{s,\,t}|\psi|^{2\sigma}\psi(s,x) ds.
\end{equation}
We denote a closed ball of $L^{q}([0,T];L^p(\bbR^d))$
\[
B_{R,\,L^{q}\left([0,\,T]\,;\,L^p\left(\bbR^d\right)\right)} : = \left\{ \psi \in L^{q}\left([0,\,T]\,;\,L^p\left(\bbR^d\right)\right); \|\psi\|_{L^{q}\left([0,\,T]\,;\,L^p\left(\bbR^d\right)\right)} \leq R\right\}.
\]
Fix $R>0$ that will be set later. For any $\psi \in B_{R,\,L^{q}([0,\,T];\,L^p(\bbR^d))}$, we apply the $L^q([0,T];L^p(\bbR^d))$ norm to~\eqref{eq:GammaMap} and deduce, thanks to~\eqref{eq:Stri1} and~\eqref{eq:Stri2},
\[
\left\|\Gamma(\psi)\right\|_{L^{q}\left([0,\,T]\,;\,L^p\left(\bbR^d\right)\right)}\leq C_1\|\psi_0\|_{L^2\left(\bbR^d\right)} + C_2|\lambda| T^{\alpha\beta}\left\|\psi\right\|_{L^{r'\left(2\sigma+1\right)}\left([0,\,T]\,;\,L^{l'\left(2\sigma+1\right)}\left(\bbR^d\right)\right)}^{2\sigma+1}.
\]
for any $(r,l)$ admissible. By choosing $(q,p) = (r,l) = \left(a,2\sigma+2\right)$, we have
\[
l' = \frac{l}{l-1} = \frac{2\sigma+2}{2\sigma +1}.
\]
Hence, we obtain, by H\"older's inequality,
\begin{multline*}
\left\|\psi\right\|_{L^{r'\left(2\sigma+1\right)}\left([0,\,T]\,;\,L^{l'\left(2\sigma+1\right)} \left(\bbR^d\right)\right)}^{2\sigma+1}\\
= \left\|\psi\right\|_{L^{r'\left(2\sigma+1\right)}\left([0,\,T]\,;\,L^{2\sigma+2}\left(\bbR^d\right)\right)}^{2\sigma+1} \leq T^{1-\frac{2}{d\sigma}} \left\|\psi\right\|_{L^{a}\left([0,\,T]\,;\,L^{2\sigma+2} \left(\bbR^d\right)\right)}^{2\sigma+1},
\end{multline*}
which gives us
\begin{multline}\label{eq:Gammaest1}
\left\|\Gamma(\psi)\right\|_{L^{q}\left([0,\,T]\,;\,L^p\left(\bbR^d\right)\right)}\\\leq C_1\|\psi_0\|_{L^2\left(\bbR^d\right)} + C_2|\lambda| T^{1+\alpha\beta-\frac{2}{d\sigma}} \left\|\psi\right\|_{L^{a}\left([0,\,T]\,;\,L^{2\sigma+2}\left(\bbR^d\right)\right)}^{2\sigma+1}.
\end{multline}
By similar computations, we obtain that, $\forall \psi_1,\psi_2\in B_{R,\,L^{q}([0,\,T]\,;\,L^p(\bbR^d))}$,
\begin{multline}\label{eq:Gammaest2}
\left\|\Gamma(\psi_1) - \Gamma(\psi_2)\right\|_{L^{q}\left([0,\,T]\,;\,L^p\left(\bbR^d\right)\right)}\\
\leq C_2|\lambda| T^{1+\alpha\beta-\frac{2}{d\sigma}} R^{2\sigma}\left\|\psi_1 - \psi_2\right\|_{L^{a}\left([0,\,T]\,;\,L^{2\sigma+2}\left(\bbR^d\right)\right)}.
\end{multline}
We remark that, since $2/d\sigma\leq 1$ and $\alpha\beta>0$, we have
\[
1+\alpha\beta-\frac{2}{d\sigma}> 0.
\]
Hence, by setting
\[
R = 2 C_1\|\psi_0\|_{L^2\left(\bbR^d\right)},
\]
and taking $T>0$ such that
\[
C_2|\lambda| T^{1+\alpha\beta-\frac{2}{d\sigma}} R^{2\sigma}< 1,
\]
we can see that $\Gamma$ is a contraction from $B_{R,\,L^{a}([0,\,T]\,;\,L^{2\sigma+2}(\bbR^d))}$ to itself. It follows from a Banach fixed point theorem that there exists a unique solution to~\eqref{eq:main}. A classical argument shows that $\psi\in\clC([0,T];L^2(\bbR^d))$. The proof of Theorem~\ref{thm:main} then follows by iterating this argument on time intervals of length $T$ since we have $\|\psi(T)\|_{L^2(\bbR^d)} = \|\psi_0\|_{L^2(\bbR^d)}$.



\begin{thebibliography}{99}
\bibitem{agrawal2001applications}Agrawal, Govind P. Applications of nonlinear fiber optics, Optics and photonics, Elsevier, 2001
\bibitem{agrawal2007nonlinear}Agrawal, Govind P. Nonlinear fiber optics, Academic Press Inc., 2007
\bibitem{belaouar2015numerical}Belaouar, Radoin; de Bouard, Anne; Debussche, Arnaud Numerical analysis of the nonlinear Schrödinger equation with white noise dispersion, Stoch. Partial Differ. Equ., Anal. Comput., Volume 3 (2015) no. 1, pp. 103-132 | Zbl
\bibitem{de1999stochastic}de Bouard, Anne; Debussche, Arnaud A Stochastic Nonlinear Schrödinger Equation with Multiplicative Noise, Commun. Math. Phys., Volume 205 (1999) no. 1, pp. 161-181 | DOI | Zbl
\bibitem{de2002effect}de Bouard, Anne; Debussche, Arnaud On the effect of a noise on the solutions of the focusing supercritical nonlinear Schrödinger equation, Probab. Theory Relat. Fields, Volume 123 (2002) no. 1, pp. 76-96 | DOI | Zbl
\bibitem{de2005blow}de Bouard, Anne; Debussche, Arnaud Blow-up for the stochastic nonlinear Schrödinger equation with multiplicative noise, Ann. Probab., Volume 33 (2005) no. 3, pp. 1078-1110 | Zbl
\bibitem{de2010nonlinear}de Bouard, Anne; Debussche, Arnaud The nonlinear Schrödinger equation with white noise dispersion, J. Funct. Anal., Volume 259 (2010) no. 5, pp. 1300-1321 | DOI | Zbl
\bibitem{catellier2016rough}Catellier, Rémi Rough linear transport equation with an irregular drift, Stoch. Partial Differ. Equ., Anal. Comput., Volume 4 (2016) no. 3, pp. 477-534 | MR | Zbl
\bibitem{cazenave2003semilinear}Cazenave, Thierry Semilinear schrödinger equations, Courant Lecture Notes in Mathematics, 10, American Mathematical Society; Courant Institute, 2003 | Zbl
\bibitem{chouk2015nonlinear}Chouk, Khalil; Gubinelli, Massimiliano Nonlinear PDEs with modulated dispersion I: Nonlinear Schrödinger equations, Commun. Partial Differ. Equations, Volume 40 (2015) no. 11, pp. 2047-2081 | DOI | Zbl
\bibitem{CATELLIER20162323}Catellier, Rémi; Gubinelli, Massimiliano Averaging along irregular curves and regularisation of ODEs, Stochastic Processes Appl., Volume 126 (2016) no. 8, pp. 2323 -2366 | MR | Zbl | DOI
\bibitem{chen2017generalized}Chen, Min; Goubet, Olivier; Mammeri, Youcef Generalized regularized long wave equation with white noise dispersion, Stoch. Partial Differ. Equ., Anal. Comput., Volume 5 (2017) no. 3, pp. 319-342 | MR | Zbl
\bibitem{da2013strong}Da Prato, Giuseppe; Flandoli, Franco; Priola, Enrico; Röckner, Michael Strong uniqueness for stochastic evolution equations in Hilbert spaces perturbed by a bounded measurable drift, Ann. Probab., Volume 41 (2013) no. 5, pp. 3306-3344 | MR | Zbl
\bibitem{debussche20111d}Debussche, Arnaud; Tsutsumi, Yoshio 1D quintic nonlinear Schrödinger equation with white noise dispersion, J. Math. Pures Appl., Volume 96 (2011) no. 4, pp. 363-376 | DOI | Zbl
\bibitem{erdHos2008quantum}Erdös, Lázló; Salmhofer, Manfred; Yau, Horng-Tzer Quantum diffusion of the random Schrödinger evolution in the scaling limit, Acta Math., Volume 200 (2008) no. 2, p. 211 | DOI | Zbl
\bibitem{erdHos2000linear}Erdös, Lázló; Yau, Horng-Tzer Linear Boltzmann equation as the weak coupling limit of a random Schrödinger equation, Commun. Pure Appl. Math., Volume 53 (2000) no. 6, pp. 667-735 | DOI | Zbl
\bibitem{fedrizzi2013noise}Fedrizzi, Ennio; Flandoli, Franco Noise prevents singularities in linear transport equations, J. Funct. Anal., Volume 264 (2013) no. 6, pp. 1329-1354 | DOI | MR | Zbl
\bibitem{flandoli2010well}Flandoli, Franco; Gubinelli, Massimiliano; Priola, Enrico Well-posedness of the transport equation by stochastic perturbation, Invent. Math., Volume 180 (2010) no. 1, pp. 1-53 | DOI | MR | Zbl
\bibitem{flandoli2011random}Flandoli, Franco Random Perturbation of PDEs and Fluid Dynamic Models : École d’été de Probabilités de Saint-Flour XL–2010, Lecture Notes in Mathematics, 2015, Springer, 2011 | DOI | Zbl
\bibitem{gess2017long}Gess, Benjamin; Souganidis, Panagiotis E. Long-Time Behavior, Invariant Measures, and Regularizing Effects for Stochastic Scalar Conservation Laws, Commun. Pure Appl. Math., Volume 70 (2017) no. 8, pp. 1562-1597 | MR | DOI | Zbl
\bibitem{hardy1928some}Hardy, Godfrey H.; Littlewood, John E. Some properties of fractional integrals. I., Math. Z., Volume 27 (1928) no. 1, pp. 565-606 | DOI | MR | Zbl
\bibitem{hardy1932some}Hardy, Godfrey H.; Littlewood, John E. Some properties of fractional integrals. II, Math. Z., Volume 34 (1931) no. 1, pp. 403-439 | DOI | MR | Zbl
\bibitem{krylov2005strong}Krylov, Nikolaĭ V.; Röckner, Michael Strong solutions of stochastic equations with singular time dependent drift, Probab. Theory Relat. Fields, Volume 131 (2005) no. 2, pp. 154-196 | DOI | MR | Zbl
\bibitem{keel1998endpoint}Keel, Markus; Tao, Terence Endpoint strichartz estimates, Am. J. Math., Volume 120 (1998) no. 5, pp. 955-980 | DOI | MR | Zbl
\bibitem{lieb2001analysis}Lieb, Elliott H.; Loss, Michael Analysis, Graduate Studies in Mathematics, 14, American Mathematical Society, 2001 | Zbl
\bibitem{marty2006splitting}Marty, Renaud On a splitting scheme for the nonlinear Schrödinger equation in a random medium, Commun. Math. Sci., Volume 4 (2006) no. 4, pp. 679-705 | DOI | MR | Zbl
\bibitem{priola2012pathwise}Priola, Enrico Pathwise uniqueness for singular SDEs driven by stable processes, Osaka J. Math., Volume 49 (2012) no. 2, pp. 421-447 | MR | Zbl
\bibitem{sobolev1938theorem}Sobolev, Sergeĭ L. On a theorem of functional analysis, Rec. Math. Moscou, n. Ser., Volume 4 (1938), pp. 471-497 | Zbl
\bibitem{veretennikov1981strong}Veretennikov, Alexander J. On strong solutions and explicit formulas for solutions of stochastic integral equations, Math. USSR, Sb., Volume 39 (1981) no. 3, pp. 387-403 | DOI | Zbl
\bibitem{xiao1997holder}Xiao, Yimin Hölder conditions for the local times and the Hausdorff measure of the level sets of Gaussian random fields, Probab. Theory Relat. Fields, Volume 109 (1997) no. 1, pp. 129-157 | DOI | Zbl
\bibitem{zvonkin1974transformation}Zvonkin, Alexander K. A transformation of the phase space of a diffusion process that removes the drift, Math. USSR, Sb., Volume 22 (1974) (1975) no. 1, pp. 129-149 | Zbl
\end{thebibliography}
\end{document}
